% !TEX root = ../main.tex
\lecture{3}{2026-10-01}{Woche 3}

\section{Daten}

\begin{definition}
    [Big Daten] Volume (Grosse Verarbeitungs, Speicherkapazität), Variety (heterogene Daten), Velocity 
\end{definition}

\subsection{Encoding}

UNICODE: Encoding (Semantik mapping binär zu Zeichen) mit maximal 32-Bit (oder $\sum_{n=0}^{32}2^n$)

\subsection{Konversion}

Bit -(:8)-> Byte -(1024)-> KByte -()


\subsection{Datenmodellierung}

\begin{definition}
    [ERM] Entity Relationship Management. Relationale Datenbanken, die strukturierte Daten in Entitäten (Tabellen) mit einander verknüpfen.
\end{definition}

\subsubsection{Grundsätze}
\begin{itemize}
    \item Atomar
    \item Eindeutig oder nicht Redundant
    \item 
\end{itemize}
\subsubsection{Beziehungen}

\begin{itemize}
    \item 1-n
    \subitem FK bei n
    \item 1-1
    \subitem entweder oder
    \item N-M 
    \subitem zwischen tabelle Lookup tabelle   
\end{itemize}

\subsubsection{SQL}

\textbf{Create Table}
\begin{verbatim}
    CREATE TABLE table_name (
        Id INTEGER PRIMARY KEY, 
        Name VARCHAR(255) NOT NULL
    );
\end{verbatim}

\textbf{Queries}

\begin{verbatim}
    SELECT * FROM table_name
    WHERE table_name.column_1 = "xxxx" AND
    table_name.column_2 = "yyyÿ:"
    ORDER BY table_name.column_2 ASC/DESC;
\end{verbatim}

\section{PageRank}
\begin{definition}
    [Page Rank]
    Innovation der Google Suchmaschine, Webseiten nach einheitlichem Kriterien bewerten und nach Relevanz vergleichen. Dämpfungsfaktor (\textit{d}) 
    \begin{equation}
        PageRank_x = (1-d) + d(\sum_{i \mid i \rightarrow x}^{x}\frac{PageRank_x}{Anzahl_{outgoinglinks_i}})
    \end{equation}
\end{definition}

\subsection{Precision/Recall}

\begin{definition}
    [Precision] Von den Gefundenen wie viele Relevanten. 
    \begin{equation}
        \frac{|R \cap G |}{|G|} = 1
    \end{equation}
    Alle Gefundenen Dokumenten (G) sind alle relevant (R). Bei < 1 wurden mehr gefunden als relevant 
\end{definition}

\begin{definition}
    [Recall] Von den Relevanten wie viele Gefunden.
      \begin{equation}
          \frac{|G \cap R |}{|R|} = 1
      \end{equation}
       Alle relevante Dokumenten wurden gefunden aber nicht unbedingt alle gefunden sind relevant.
\end{definition}



\subsection{Aufgabentypen}
\begin{itemize}
    \item Kovertieren Bin -> Dec  oder Dec -> Bin
    \item KOnvertieren Bin -> ASCII oder ASCII->Bin
    \item Speicherplatz eines Dokuments berechnen anhand UNICODE, Seiten, Zeilen, Zeichen
    \item PageRank eines netzwerks mit neuem Dokument bewerten
    \item Recall udn Precision einer Suchmaschine Bewerten
\end{itemize}